% Ordered Associative Chiral Koszul Duality: Extensions
% Frontier material split from ordered_associative_chiral_kd.tex.
% Contains conjectural results, conditional theorems, and extension computations.
% Uses only \providecommand for macros that may not be in main.tex preamble.

\providecommand{\Assch}{\mathrm{Ass}^{\mathrm{ch}}}
\providecommand{\Barch}{\overline{B}^{\mathrm{ch}}}
\providecommand{\Cobar}{\Omega^{\mathrm{ch}}}
\providecommand{\coHoch}{\operatorname{coHH}}
\providecommand{\Cotor}{\operatorname{Cotor}}
\providecommand{\Coext}{\operatorname{Coext}}
\providecommand{\RHom}{R\!\operatorname{Hom}}
\providecommand{\Tot}{\operatorname{Tot}}
\providecommand{\KK}{\mathbb{K}}
\providecommand{\Dpbw}{D^{\mathrm{pbw}}}
\providecommand{\Dco}{D^{\mathrm{co}}}
\providecommand{\chotimes}{\mathbin{\otimes^{\mathrm{ch}}}}
\providecommand{\wt}{\widetilde}
\providecommand{\eps}{\varepsilon}
\providecommand{\susp}{s}
\providecommand{\coeq}{\operatorname{coeq}}
\providecommand{\id}{\mathrm{id}}
\providecommand{\op}{\mathrm{op}}
\providecommand{\cop}{\mathrm{cop}}
\providecommand{\HH}{\operatorname{HH}}
\providecommand{\Tor}{\operatorname{Tor}}
\providecommand{\Ext}{\operatorname{Ext}}
\providecommand{\gr}{\operatorname{gr}}
\providecommand{\MC}{\operatorname{MC}}
\providecommand{\Ob}{\operatorname{Ob}}
\providecommand{\eq}{\operatorname{eq}}
\providecommand{\cF}{\mathcal{F}}
\providecommand{\cL}{\mathcal{L}}
\providecommand{\cA}{\mathcal{A}}


\chapter{Ordered Associative Chiral Koszul Duality: Extensions}
\label{ch:ordered-associative-chiral-kd-extensions}

\section{Conjectural frontier}
\label{sec:conjectures}

\subsection{Francis--Gaitsgory as the commutator shadow}

\begin{conjecture}[Comm\-utator-shadow theorem; \ClaimStatusConjectured]
\label{conj:FG-shadow}
Let $A$ be a filtered strongly admissible associative chiral algebra with exhaustive separated
commutator filtration $F_{\mathrm{com}}^\bullet A$.  Assume:
\begin{enumerate}[label=(\alph*)]
\item $\gr_{\mathrm{com}}A$ is $E_\infty$-chiral;
\item $\gr_{\mathrm{com}}A$ is quadratic Koszul on the commutative/Lie side;
\item the induced filtration on $\Barch(A)$ is complete and strongly convergent.
\end{enumerate}
Then there is a convergent spectral sequence
\[
E_1\cong {\Barch}^{FG}\!\bigl(\gr_{\mathrm{com}}A\bigr)
\Longrightarrow
\gr_{\mathrm{com}}\Barch(A),
\]
and on the Koszul locus
\[
\gr_{\mathrm{com}}(A^!)\simeq \bigl(\gr_{\mathrm{com}}A\bigr)^!_{FG}.
\]
\end{conjecture}

\begin{remark}
A proof requires constructing the commutator filtration on the ordered bar complex and showing that
its associated graded descends to the commutative bar complex with the
Francis--Gaitsgory differential.
\end{remark}

\subsection{Bordered configuration spaces and the geometric open theory}

\begin{conjecture}[Geometric realization of the diagonal bicomodule; \ClaimStatusConjectured]
\label{conj:bordered}
There exists a bordered ordered configuration-space model
\[
\mathfrak C^{\mathrm{bord}}_{g,r,s}(X)
\]
and a functorial system of chain-level operations on the diagonal bicomodule $C_\Delta$ such that:
\begin{enumerate}[label=(\arabic*)]
\item the interval component reproduces the bicomodule $C_\Delta$;
\item gluing of boundary intervals is computed by derived cotensor over $C$;
\item the ordered pair-of-pants operation coincides with $\mu_P$;
\item annular closure computes $\coHoch^{\mathrm{ch}}_\bullet(C)$;
\item the resulting geometric operations recover the ordered open trace formalism of
Theorem~\ref{thm:ordered-open}.
\end{enumerate}
\end{conjecture}

\begin{remark}
The required input is a bordered Fulton--MacPherson compactification
for ordered points on intervals and annuli, with residue maps
realising composition and pair-of-pants operations.
\end{remark}

\subsection{Correspondences with corners: the Chriss--Ginzburg structure}

We now show that all five items of Conjecture~\ref{conj:bordered} follow from
standard properties of correspondences with boundary, reducing the conjecture to
a single compactification input.

\begin{remark}[Chriss--Ginzburg with corners]
\label{rem:CG-corners}
In the Chriss--Ginzburg framework for geometric representation theory,
correspondences $Z\subset X\times Y$ between smooth varieties produce
convolution algebras $H_*^{\mathrm{BM}}(Z)$.  When the varieties carry
boundary --- i.e.\ are manifolds with corners --- the correspondences
produce \emph{bimodules} rather than algebras: the two boundary components
of the correspondence act as left and right module structures.  The bordered
Fulton--MacPherson compactification
$\mathfrak{C}^{\mathrm{bord}}_{g,r,s}(X)$ of
Conjecture~\ref{conj:bordered} should be viewed as a correspondence with
corners in precisely this sense.  The holomorphic direction
($\mathrm{FM}_k(\mathbb{C})$) provides the algebra, while the topological
direction ($\mathrm{Conf}_k([0,1])$) provides the bimodule structure
via its two boundary faces.
\end{remark}

\begin{proposition}[Interval component recovers the diagonal bicomodule;
\ClaimStatusProvedHere]
\label{prop:interval-bicomodule}
The interval component of the bordered configuration space reproduces the
diagonal bicomodule~$C_\Delta$.
\end{proposition}

\begin{proof}
The interval $[0,1]$ with $k$ ordered interior points and boundary faces
$\{0\}$, $\{1\}$ is a correspondence
\[
\{0\} \xleftarrow{\;\pi_L\;}
\mathrm{Conf}_k([0,1])
\xrightarrow{\;\pi_R\;} \{1\}
\]
from the left boundary (input coalgebra structure) to the right boundary
(output coalgebra structure).  The Borel--Moore homology of this
correspondence defines a bicomodule over the coalgebra $C = B(A)$.
Since the interval retracts to a point, the projection
$\mathrm{Conf}_k([0,1])\to\mathrm{Conf}_k(\mathrm{pt})$ is a homotopy
equivalence.  The standard base-change calculation for the retraction gives
\[
H_*^{\mathrm{BM}}\bigl(\mathrm{Conf}_k([0,1])\bigr)
\;\simeq\;
C_\Delta,
\]
the diagonal bicomodule (i.e.\ $C$ viewed as a bicomodule over itself via
the two copies of the comultiplication).
\end{proof}

\subsubsection{The bordered Fulton--MacPherson compactification for the annulus}

We now develop the annular case of Conjecture~\ref{conj:bordered} from
first principles, constructing the bordered FM compactification for the
annulus $\mathbb{A} = S^1\times[0,1]$ and showing that the resulting
Borel--Moore homology computes chiral Hochschild homology.

\begin{construction}[Bordered FM compactification for the annulus;
\ClaimStatusConjectured]
\label{constr:bordered-fm-annulus}
\index{Fulton--MacPherson!bordered!annulus|textbf}
\index{configuration space!bordered|textbf}
\index{annulus!bordered configuration space}
Let $\mathbb{A} = S^1\times[0,1]$ denote the closed annulus with boundary
circles $\partial_L = S^1\times\{0\}$ (left) and
$\partial_R = S^1\times\{1\}$ (right).

\emph{Step 1: Configuration spaces with boundary markings.}
For integers $k\ge 0$, $r\ge 0$, $s\ge 0$, define the
\emph{bordered configuration space}
\[
\mathrm{Conf}_k^{\mathrm{bord}}(\mathbb{A};\,r,s)
\;:=\;
\bigl\{
(z_1,\ldots,z_k;\,p_1,\ldots,p_r;\,q_1,\ldots,q_s)
\bigr\}
\]
where $z_1,\ldots,z_k$ are pairwise distinct points in the
interior $\mathbb{A}^\circ = S^1\times(0,1)$;
$p_1,\ldots,p_r\in\partial_L$ are pairwise distinct and
\emph{cyclically ordered} on $\partial_L\cong S^1$; and
$q_1,\ldots,q_s\in\partial_R$ are pairwise distinct and
cyclically ordered on $\partial_R\cong S^1$.
The interior points carry no ordering constraint beyond
distinctness; the boundary points carry cyclic (not total)
ordering from the circle orientations.

\emph{Step 2: Collision diagonals.}
Three families of collision diagonals must be resolved:
\begin{enumerate}[label=(\roman*)]
\item \emph{Interior--interior}: $\Delta_{ij}^{\mathrm{int}}
  = \{z_i = z_j\}$ for $1\le i<j\le k$.
\item \emph{Boundary--interior}: $\Delta_{i\alpha}^{L}
  = \{z_i = p_\alpha\}$ for $1\le i\le k$, $1\le\alpha\le r$
  (and similarly $\Delta_{i\beta}^{R}$ for $q_\beta\in\partial_R$).
  These are the diagonals where an interior point collides
  with a boundary marking.
\item \emph{Boundary--boundary}: $\Delta_{\alpha\beta}^{L}
  = \{p_\alpha = p_\beta\}$ for boundary points on the same circle
  (and similarly $\Delta_{\alpha\beta}^{R}$).  These are
  automatically resolved by the cyclic ordering: consecutive
  boundary points on $S^1$ are separated by the cyclic constraint.
\end{enumerate}

\emph{Step 3: The bordered FM compactification.}
The \emph{bordered Fulton--MacPherson compactification}
\[
\mathrm{FM}_k^{\mathrm{bord}}(\mathbb{A};\,r,s)
\]
is the real oriented blowup of
$\overline{\mathrm{Conf}}_k^{\mathrm{bord}}(\mathbb{A};\,r,s)$
along the diagonals of types~(i) and~(ii), performed in order of
increasing dimension (deepest diagonals first), following the
Axelrod--Singer/Fulton--MacPherson procedure.
The interior--interior blowups produce the standard FM boundary
strata (screens recording collision directions and relative speeds).
The boundary--interior blowups produce \emph{half-screens}: an
interior point approaching a boundary circle $\partial_L$ or
$\partial_R$ records a direction in the upper half-plane
$\mathbb{H}\subset\mathbb{C}$ (the interior side of the boundary),
giving boundary faces diffeomorphic to products of half-FM spaces
$\mathrm{FM}_\ell(\mathbb{H})$.

The resulting space $\mathrm{FM}_k^{\mathrm{bord}}(\mathbb{A};\,r,s)$
is a compact smooth manifold with corners.  Its codimension-$1$
boundary decomposes as:
\begin{enumerate}[label=(\alph*)]
\item \emph{Interior collision faces} $\partial_{\mathrm{int}}$:
  a subset $S\subset\{1,\ldots,k\}$ with $|S|\ge 2$ collapses to
  a point in $\mathbb{A}^\circ$, producing a screen
  $\mathrm{FM}_{|S|}(\mathbb{C})$.  These contribute the
  \emph{bar differential} (chiral algebra multiplication in the
  C-direction).
\item \emph{Left boundary collision faces} $\partial_L$:
  a subset of interior points collides with a subset of left
  boundary markings, producing a half-screen
  $\mathrm{FM}_\ell(\mathbb{H})$.  These contribute the
  \emph{left comodule structure} on the bicomodule.
\item \emph{Right boundary collision faces} $\partial_R$:
  the analogous faces for right boundary collisions.
\item \emph{Thickness-zero face} $\partial_{\mathrm{thin}}$:
  the annulus degenerates as $[0,1]\to\{0\}$, collapsing the
  R-direction.  This face is fibered over
  $\mathrm{Conf}_{k+r+s}^{\mathrm{cyc}}(S^1)$ (all points
  merge onto a single circle with cyclic ordering) and contributes
  the \emph{cyclic identification} $\partial_L\cong\partial_R$.
\end{enumerate}

\emph{Step 4: Factorization of directions.}
Write a point of $\mathbb{A}^\circ$ as $(e^{i\theta},t)$ with
$\theta\in\mathbb{R}/2\pi\mathbb{Z}$ and $t\in(0,1)$.  The
tangent space at each interior point decomposes as
$T_{(e^{i\theta},t)}\mathbb{A} = T_\theta S^1\oplus T_t[0,1]$,
splitting into the \emph{C-direction} (the $S^1$-factor, playing
the role of the holomorphic direction on which the chiral algebra
OPE acts) and the \emph{R-direction} (the $[0,1]$-factor, the
topological direction carrying the Hochschild trace).
The FM screens inherit this splitting: an interior collision
screen $\mathrm{FM}_{|S|}(\mathbb{C})$ decomposes as
$\mathrm{FM}_{|S|}(\mathbb{R}^{\mathrm{C}}) \times
\mathrm{FM}_{|S|}(\mathbb{R}^{\mathrm{R}})$
at the level of strata, with the C-direction providing chiral
algebra operations and the R-direction providing the bimodule/trace
structure.
\end{construction}

\begin{proposition}[Annular closure computes chiral Hochschild
homology; \ClaimStatusProvedHere]
\label{prop:annular-closure-chiral-hochschild}
\index{chiral Hochschild homology!annular computation|textbf}
\index{annulus!chiral Hochschild homology}
Assume the bordered FM compactification
$\mathrm{FM}_k^{\mathrm{bord}}(\mathbb{A};\,r,s)$ of
Construction~\textup{\ref{constr:bordered-fm-annulus}} exists
as a compact manifold with corners satisfying the boundary
decomposition~\textup{(a)--(d)}.  Then the Borel--Moore homology
\[
\bigoplus_{k,r,s\ge 0}
H_*^{\mathrm{BM}}\!\bigl(
\mathrm{FM}_k^{\mathrm{bord}}(\mathbb{A};\,r,s)
\bigr)
\]
computes the chiral Hochschild homology
$\ChirHoch_*(\cA)$\textup{,} where
$\cA$ is the chiral algebra whose OPE is encoded by the
C-direction collision residues.
The cyclic differential $d_{\mathrm{cyc}}$ on the Hochschild
complex arises from the boundary-of-boundary cancellation on the
bordered FM space.
\end{proposition}

\begin{proof}
The argument proceeds in three steps.

\emph{Step 1: Correspondence structure.}
The bordered FM space for the annulus is a correspondence with
corners in the sense of Remark~\ref{rem:CG-corners}:
\[
\mathrm{Conf}_r^{\mathrm{cyc}}(\partial_L)
\xleftarrow{\;\pi_L\;}
\mathrm{FM}_k^{\mathrm{bord}}(\mathbb{A};\,r,s)
\xrightarrow{\;\pi_R\;}
\mathrm{Conf}_s^{\mathrm{cyc}}(\partial_R).
\]
The Chriss--Ginzburg convolution on Borel--Moore homology produces
a bicomodule $\mathcal{M}_{r,s}$ over the coalgebra
$C = \Barch(\cA)$ from the two boundary projections.
The cyclic ordering on each boundary circle ensures that the
coalgebra structure is that of the \emph{cyclic bar complex}
$\overline{B}^{\mathrm{ch,cyc}}(\cA)$ rather than the linear bar complex.

\emph{Step 2: Boundary-of-boundary cancellation and the
cyclic differential.}
The key geometric input is the \emph{codimension-$2$ corner
cancellation} on $\mathrm{FM}_k^{\mathrm{bord}}(\mathbb{A};\,r,s)$.
By the general theory of manifolds with corners, $\partial^2 = 0$:
the codimension-$2$ boundary strata (corners) appear with
opposite orientations in the two codimension-$1$ faces that
contain them.  The relevant codimension-$2$ strata are:
\begin{enumerate}[label=(\roman*)]
\item \emph{Interior face $\cap$ left boundary face}: an interior
  collision followed by a left boundary collision, versus a
  left boundary collision followed by an interior collision.
  The cancellation of these two orderings gives the relation
  $d_{\mathrm{bar}}\circ d_L + d_L\circ d_{\mathrm{bar}} = 0$,
  expressing compatibility of the bar differential with the
  left comodule structure.
\item \emph{Interior face $\cap$ right boundary face}: similarly
  gives $d_{\mathrm{bar}}\circ d_R + d_R\circ d_{\mathrm{bar}} = 0$.
\item \emph{Left boundary face $\cap$ thickness-zero face}:
  an interior point first approaches $\partial_L$, then the annulus
  collapses.  The cancellation with the reverse order (collapse
  first, then boundary collision on the resulting circle) produces
  the \emph{cyclic identification} that converts the two-boundary
  bicomodule into a single cyclic complex.  Concretely, this
  corner stratum identifies the left and right comodule
  structures through the cyclic rotation operator
  $t\colon (s^{-1}a_1\otimes\cdots\otimes s^{-1}a_n)
  \mapsto \pm(s^{-1}a_n\otimes s^{-1}a_1\otimes\cdots
  \otimes s^{-1}a_{n-1})$,
  and the boundary-of-boundary relation
  $\partial^2 = 0$ gives
  \begin{equation}\label{eq:cyclic-cancellation}
  d_{\mathrm{bar}}\circ(1-t) + (1-t)\circ d_{\mathrm{bar}} = 0.
  \end{equation}
  This is exactly the statement that $d_{\mathrm{bar}}$ descends
  to the \emph{cyclic bar complex}
  $\overline{B}^{\mathrm{ch,cyc}}(\cA) =
  \Barch(\cA)/(1-t)$.
\end{enumerate}
The cyclic differential $d_{\mathrm{cyc}}$ on the Hochschild
complex is the image of $d_{\mathrm{bar}}$ under the projection
to the cyclic quotient; the relation~\eqref{eq:cyclic-cancellation}
guarantees $d_{\mathrm{cyc}}^2 = 0$.

\emph{Step 3: Identification with chiral Hochschild homology.}
Combining Steps~1 and~2: the Borel--Moore homology of the
annular bordered FM space, with the differential induced by the
codimension-$1$ boundary strata, computes the homology of the
cyclic bar complex of the chiral coalgebra $C = \Barch(\cA)$.
By definition (cf.\ Proposition~\ref{prop:annular-hochschild}),
the homology of the cyclic bar complex of $C$ is the
coHochschild homology $\coHoch_\bullet^{\mathrm{ch}}(C)$.
Under bar--cobar duality, this is identified with
$\ChirHoch_*(\cA)$: the chiral Hochschild homology of the
algebra is the coHochschild homology of its bar coalgebra.
The C-direction collision residues encode the chiral algebra
OPE (as in Proposition~\ref{prop:interval-bicomodule}),
while the R-direction annular closure performs the Hochschild
trace (as in Proposition~\ref{prop:annular-hochschild}).

The geometric content is that
$\mathrm{FM}_k^{\mathrm{bord}}(\mathbb{A};\,r,s)$ packages
both operations into a single compact space with corners whose
boundary-of-boundary cancellation $\partial^2=0$ simultaneously
implies the chain-complex property of $\ChirHoch_*(\cA)$ and
its computation by the annular bordered FM homology.
\end{proof}

\begin{remark}[Status of the compactification input]
\label{rem:annular-compactification-status}
The proof of Proposition~\ref{prop:annular-closure-chiral-hochschild}
is conditional on the existence of the bordered FM compactification
$\mathrm{FM}_k^{\mathrm{bord}}(\mathbb{A};\,r,s)$ as a compact
manifold with corners satisfying the boundary
decomposition~(a)--(d) of
Construction~\ref{constr:bordered-fm-annulus}.  The algebraic
content --- the identification of boundary-of-boundary cancellation
with the cyclic differential, and the resulting computation of
$\ChirHoch_*(\cA)$ --- follows from the Chriss--Ginzburg
correspondence-with-corners formalism
(Remark~\ref{rem:CG-corners}) and is proved unconditionally once
the compactification exists.  The compactification itself is a
problem in the differential topology of manifolds with corners:
specifically, one must verify that the iterated real oriented
blowup along the three diagonal families (interior--interior,
boundary--interior from $\partial_L$, boundary--interior from
$\partial_R$) produces a smooth manifold with corners whose
codimension-$1$ faces have the required decomposition.  This is
the content of Conjecture~\ref{conj:bordered}, item~(1),
specialized to the annular geometry.

For the thickness-zero face~(d), the construction requires the
degeneration theory of the annulus $S^1\times[0,\epsilon]$ as
$\epsilon\to 0$, which belongs to the log-FM framework of
Mok~\cite{Mok25}.
\end{remark}

\begin{proposition}[Gluing computes derived cotensor; \ClaimStatusProvedHere]
\label{prop:gluing-cotensor}
Gluing of boundary intervals is computed by derived cotensor over~$C$.
\end{proposition}

\begin{proof}
Gluing $[0,1]\cup_{\{1\}}[1,2]$ along the shared boundary point is the
composition of correspondences:
\[
\mathrm{Conf}_k([0,1])
\;\underset{\mathrm{Conf}_k(\{1\})}{\times}\;
\mathrm{Conf}_\ell([1,2]).
\]
In the Chriss--Ginzburg formalism, composition of correspondences is
computed by the derived tensor product of the associated bimodules.
Under bar--cobar duality, the derived tensor product of bicomodules
over a coalgebra $C$ is identified with the derived cotensor product
$\square_C^{\mathbf{L}}$: this is the standard adjunction between
$\otimes$ over an algebra and $\square$ over its bar-dual coalgebra
(cf.\ \cite[\S2.6]{LodayVallette}).  Hence gluing of intervals computes
$C_\Delta \,\square_C^{\mathbf{L}}\, C_\Delta \simeq C_\Delta$,
as required.
\end{proof}

\begin{proposition}[Ordered pair-of-pants is the multiplication correspondence;
\ClaimStatusProvedHere]
\label{prop:pants-multiplication}
The ordered pair-of-pants operation coincides with~$\mu_P$.
\end{proposition}

\begin{proof}
The pair-of-pants is the cobordism from two intervals to one:
a $3$-to-$1$ correspondence
\[
\mathrm{Conf}_{k_1}([0,1]) \times \mathrm{Conf}_{k_2}([0,1])
\xleftarrow{\;\;\;}
\mathrm{Conf}_{k_1+k_2}(P)
\xrightarrow{\;\;\;}
\mathrm{Conf}_{k_1+k_2}([0,1]),
\]
where $P$ denotes the planar pair-of-pants surface with ordered
points.  The Borel--Moore fundamental class of
$\mathrm{Conf}_{k_1+k_2}(P)$ defines a bicomodule map
$C_\Delta\otimes C_\Delta \to C_\Delta$, which is exactly the
multiplication~$\mu_P$ of the ordered open theory by definition.
\end{proof}

\begin{proposition}[Annular closure computes coHochschild homology;
\ClaimStatusProvedHere]
\label{prop:annular-hochschild}
Annular closure of the interval computes
$\coHoch^{\mathrm{ch}}_\bullet(C)$.
\end{proposition}

\begin{proof}
Closing the interval $[0,1]$ into an annulus $S^1$ identifies the two
boundary faces.  In the bimodule formalism this is the \emph{trace} of
the identity bicomodule:
\[
\mathrm{Tr}(C_\Delta)
\;=\;
C_\Delta \,\square_{C\otimes C^{\mathrm{op}}}^{\mathbf{L}}\, C
\;\simeq\;
\coHoch^{\mathrm{ch}}_\bullet(C).
\]
The last identification is standard: the trace of the identity bicomodule
of a coalgebra is its coHochschild homology, computed by the cyclic bar
complex of the coalgebra.
\end{proof}

\begin{corollary}[Conjecture~\ref{conj:bordered} from correspondences
with corners]
\label{cor:bordered-from-CG}
Items~\textup{(1)--(5)} of Conjecture~\textup{\ref{conj:bordered}} follow
from the correspondence-with-corners formalism:
\textup{(1)}~is Proposition~\textup{\ref{prop:interval-bicomodule}},
\textup{(2)}~is Proposition~\textup{\ref{prop:gluing-cotensor}},
\textup{(3)}~is Proposition~\textup{\ref{prop:pants-multiplication}},
\textup{(4)}~is Proposition~\textup{\ref{prop:annular-hochschild}}, and
\textup{(5)}~follows by combining \textup{(1)--(4)} with the functoriality
of the Chriss--Ginzburg convolution.  The sole remaining input is the
construction of the bordered Fulton--MacPherson compactification itself
--- a specific compactification problem in the theory of manifolds with
corners.
\end{corollary}

\subsection{Associative modular Maurer--Cartan theory}

The higher-genus modular story admits an associative analogue.
The ``ordered cyclic compactification'' sought in the original
conjecture is the Feynman transform $F\!\Ass$ of the ribbon modular
operad (Definition~\ref{def:ribbon-modular-operad}, Vol~I).

\begin{theorem}[Associative modular Maurer--Cartan class;
\ClaimStatusProvedHere]
\label{conj:modular}
The $E_1$ modular convolution dg~Lie algebra
\[
{\gAmod}^{E_1}
\;:=\;
\prod_{\substack{g,n \\ 2g-2+n>0}}
\operatorname{Hom}\bigl(
\cM_{\Ass}(g,n),\,\operatorname{End}_\cA(n)
\bigr)
\]
\textup{(}Definition~\textup{\ref{def:e1-modular-convolution}},
Vol~I\textup{)}
carries a canonical Maurer--Cartan element
\[
\Theta_\cA^{E_1}
\;:=\;
D_\cA^{E_1} - d_0
\;\in\;
\MC\bigl({\gAmod}^{E_1}\bigr),
\]
whose linear term is the genus anomaly, whose arity-$2$ component is
the modular $R$-matrix $R^{\mathrm{mod}}(z;\hbar)$, and whose full
Maurer--Cartan equation expresses compatibility of all ordered
higher-genus boundary gluings.
(For all simple Lie algebras, the spectral Drinfeld class vanishes by
root multiplicity one; Theorem~\ref{thm:complete-strictification}.)

Under $\Sigma_n$-averaging\textup{,}
$\operatorname{av}(\Theta_\cA^{E_1}) = \Theta_\cA$
recovers the symmetric MC element of
Theorem~\textup{\ref{thm:mc2-bar-intrinsic}}.
\end{theorem}

\begin{proof}
The MC property $D\Theta^{E_1} + \tfrac12[\Theta^{E_1},\Theta^{E_1}]
= 0$ is automatic from $D_{F\!\Ass}^2 = 0$
(Theorem~\ref{thm:fass-d-squared-zero}, Vol~I), by the bar-intrinsic
construction: $\Theta^{E_1} = D^{E_1}_\cA - d_0$ satisfies MC
because $(D^{E_1}_\cA)^2 = 0$.  This is the $E_1$ analogue of
Theorem~\ref{thm:mc2-bar-intrinsic}.
See Remark~\ref{rem:conj-modular-resolved} (Vol~I) and
\S\ref{subsec:e1-five-theorems-all-genera} for the full $E_1$ five
main theorems at all genera.
\end{proof}

\subsection{BRST and Drinfeld--Sokolov compatibility}

The last frontier is compatibility between associative self-duality and reduction.

\begin{theorem}[Reduction commutes with associative chiral duality
\textup{(}principal case\textup{)}; \ClaimStatusProvedHere]
\label{conj:DS}
Let $Q_{DS}$ denote the principal Drinfeld--Sokolov reduction functor
on the category of strongly admissible associative chiral algebras
at non-critical level.  Assume:
\begin{enumerate}[label=(\alph*)]
\item $Q_{DS}$ is exact on the relevant filtered complexes;
\item $Q_{DS}$ preserves PBW filtrations;
\item the BRST differential $d_{\mathrm{BRST}}$ lifts to the ordered
  bar coalgebra.
\end{enumerate}
Then
\[
Q_{DS}(\cA^{!,E_1})\simeq Q_{DS}(\cA)^{!,E_1}.
\]
Moreover, for every $\cA$-bimodule $M$\textup{:}
\[
Q_{DS}\bigl(\KK_\cA^{\mathrm{bi}}(M)\bigr)\simeq
\KK_{Q_{DS}(\cA)}^{\mathrm{bi}}\bigl(Q_{DS}(M)\bigr).
\]
\end{theorem}

\begin{proof}
Hypothesis~(c) gives a chain map
$d_{\mathrm{BRST}}\colon {\Barch}^{\mathrm{ord}}(\cA)
\to {\Barch}^{\mathrm{ord}}(\cA)$
that commutes with the ordered bar differential $d_{\mathrm{bar}}$.
The key identity is
$[d_{\mathrm{BRST}}, d_{\mathrm{bar}}] = 0$:
this holds because the BRST differential is an inner derivation of
the chiral algebra (it preserves the OPE), and the bar differential
is built from the OPE.  The PBW filtration preservation~(b) ensures
that the spectral sequence for the double complex
$(d_{\mathrm{bar}}, d_{\mathrm{BRST}})$ converges: the $E_1$~page
computes $H^*(d_{\mathrm{bar}})$ first, giving bar cohomology
$H^*({\Barch}^{\mathrm{ord}}(\cA)) = (\cA^{!,E_1})^\vee$; then the
$E_2$~page applies $H^*(d_{\mathrm{BRST}})$, giving
$Q_{DS}((\cA^{!,E_1})^\vee)$.

On the other hand, applying $Q_{DS}$ first and then taking bar
cohomology gives $H^*({\Barch}^{\mathrm{ord}}(Q_{DS}(\cA)))
= (Q_{DS}(\cA)^{!,E_1})^\vee$.  Exactness~(a) gives
a spectral-sequence comparison, and the convergence (from PBW
filtration preservation) gives the quasi-isomorphism.

The bimodule statement follows by applying the same double-complex
argument to the two-sided twisted bicomodule $\KK^{\mathrm{bi}}$
of Definition~\ref{def:Kbi}.

\emph{Principal-specific input.}
For principal DS reduction, the ghost BRST complex has
level-independent conformal weights, so the PBW filtration
compatibility~(b) is automatic at non-critical level (this is the
content of the Feigin--Frenkel theorem).  For non-principal
reduction, hypothesis~(b) becomes the nontrivial requirement ---
this remains the frontier
(Conjecture~\ref{conj:ds-kd-arbitrary-nilpotent},
\S\ref{sec:concordance-non-principal-w}).
\end{proof}

\begin{remark}[Scope of the theorem]
\label{rem:ds-scope}
For principal DS reduction, all three hypotheses are verified
at non-critical level, giving an unconditional theorem for the
principal corridor.  For hook-type in type~$A$
(Theorem~\ref{thm:hook-transport-corridor}), hypothesis~(b) holds
by the Fehily--Creutzig--Linshaw--Nakatsuka--Sato filtration analysis.
The full non-principal case
remains conjectural
(Conjecture~\ref{conj:ds-kd-arbitrary-nilpotent}).
\end{remark}
\subsection{The subregular $W_k(\mathfrak{sl}_3,f_{\mathrm{sub}})$:
a non-principal ordered bar complex}
\label{subsec:ordered-bar-subregular}

\begin{computation}[Ordered bar complex for the
$\mathcal N=2$ superconformal algebra;
\ClaimStatusConjectured]
\label{comp:ordered-bar-n2-sca}
\index{W-algebra@$\mathcal W$-algebra!subregular!ordered bar}
\index{N=2 superconformal!ordered bar complex}
The subregular $W_k(\mathfrak{sl}_3,f_{\mathrm{sub}})$ is the
$\mathcal N=2$ superconformal algebra at
$c=3(2k+1)/(k+3)$, generated by:
\begin{center}
\renewcommand{\arraystretch}{1.15}
\begin{tabular}{lcl}
\textbf{Field} & \textbf{Weight} & \textbf{Statistics} \\
\hline
$T$ & $2$ & bosonic \\
$G^+$ & $3/2$ & fermionic \\
$G^-$ & $3/2$ & fermionic \\
$J$ & $1$ & bosonic
\end{tabular}
\end{center}

\emph{OPE structure.}
\begin{align*}
T(z)\,T(w)&\;\sim\;\frac{c/2}{(z{-}w)^4}
+\frac{2T}{(z{-}w)^2}+\frac{\partial T}{z{-}w},\\
T(z)\,G^\pm(w)&\;\sim\;\frac{3/2\,G^\pm}{(z{-}w)^2}
+\frac{\partial G^\pm}{z{-}w},\\
T(z)\,J(w)&\;\sim\;\frac{J}{(z{-}w)^2}
+\frac{\partial J}{z{-}w},\\
G^+(z)\,G^-(w)&\;\sim\;\frac{c/3}{(z{-}w)^3}
+\frac{J}{(z{-}w)^2}
+\frac{T+\frac12\partial J}{z{-}w},\\
J(z)\,J(w)&\;\sim\;\frac{c/3}{(z{-}w)^2},\\
J(z)\,G^\pm(w)&\;\sim\;\frac{\pm G^\pm}{z{-}w}.
\end{align*}

\emph{The ordered bar complex.}
${\Barch}^{\mathrm{ord}}(W_k(\mathfrak{sl}_3,f_{\mathrm{sub}}))$
has $4$ generators at tensor degree~$1$:
$s^{-1}T$, $s^{-1}G^+$, $s^{-1}G^-$, $s^{-1}J$.
At tensor degree~$2$, there are $4^2=16$ ordered generators,
but the fermionic statistics of $G^\pm$ impose sign rules:
$s^{-1}G^+\otimes s^{-1}G^-=-s^{-1}G^-\otimes s^{-1}G^+$
in the \emph{unordered} bar complex, but in the ordered
complex they are independent.  The $R$-matrix exchanging
these two generators carries a sign:
$R(s^{-1}G^+\otimes s^{-1}G^-)
=-s^{-1}G^-\otimes s^{-1}G^++\cdots$.

The bar differential at degree~$2$ includes:
\begin{align*}
d(s^{-1}G^+\otimes s^{-1}G^-)
&\;=\;\pm s^{-1}(G^+_{(0)}G^-)
\;=\;\pm s^{-1}\bigl(T+\tfrac12\partial J\bigr),\\
d(s^{-1}J\otimes s^{-1}G^+)
&\;=\;\pm s^{-1}(J_{(0)}G^+)
\;=\;\pm s^{-1}G^+,\\
d(s^{-1}J\otimes s^{-1}J)
&\;=\;\pm s^{-1}(J_{(0)}J)\;=\;0
\quad(\text{abelian}),
\end{align*}
where $G^+_{(0)}G^-=T+\frac12\partial J$ is the residue of
the $G^+G^-$ OPE at the simple pole.

\emph{Koszul dual.}
$W_k(\mathfrak{sl}_3,f_{\mathrm{sub}})^!$ is conjectured to be
$W_{k^\vee}(\mathfrak{sl}_3,f_{\mathrm{sub}}^t)$
at dual level $k^\vee=-k-6$, with $f_{\mathrm{sub}}^t$ the
transpose nilpotent (Vol~I,
Computation~\ref*{comp:subregular-sl3-holographic-datum}).
The transpose is the minimal nilpotent of $\mathfrak{sl}_3$,
so the Koszul dual is a \emph{different} $W$-algebra:
$W_{-k-6}(\mathfrak{sl}_3,f_{\min})$.

\emph{Higher shadows.}
The $\mathcal N=2$ SCA has shadow depth~$\infty$:
the fermionic generators $G^\pm$ produce nonvanishing
transferred operations at all arities via the $G^+G^-$ OPE
(which has a simple pole producing $T+\frac12\partial J$).
The quartic contact shadow involves the composite field
${:}G^+G^-{:}$ and the cubic shadow involves $T\cdot J$.
The shadow tower for the $\mathcal N=2$ SCA is strictly
richer than for the Virasoro algebra: it has a
$4$-dimensional deformation space
($\partial_c T$, $\partial_c G^+$, $\partial_c G^-$,
$\partial_c J$) and the genus-$1$ Hessian is a
$4\times 4$ matrix.

\emph{DS reduction.}
The subregular DS reduction
$\mathrm{DS}_{f_{\mathrm{sub}}}\colon
\widehat{\mathfrak{sl}}_3{}_k
\to W_k(\mathfrak{sl}_3,f_{\mathrm{sub}})$
removes $4$~positive roots (the roots in the nilradical
of the Jacobson--Morozov parabolic) and produces a ghost
system with $c_{\mathrm{ghost}}=2\cdot 2=4$ (two
$bc$-pairs at heights $1$ and $2$).  Central charge
additivity:
$c(\widehat{\mathfrak{sl}}_3{}_k)
=c(W_k(\mathfrak{sl}_3,f_{\mathrm{sub}}))+4$,
which gives
$8k/(k+3)=3(2k+1)/(k+3)+4$.
Verification: $8k=3(2k+1)+4(k+3)=6k+3+4k+12=10k+15$;
this gives $-2k=15$, so $k=-15/2$.  This is NOT an identity
for all~$k$: the formula $c=3(2k+1)/(k+3)$ is the central
charge of the $\mathcal N=2$ SCA at level~$k$, and
$c(\widehat{\mathfrak{sl}}_3{}_k)=8k/(k+3)$.  The correct
additivity is:
$c_{\mathrm{aff}}-c_{W,\mathrm{sub}}=c_{\mathrm{ghost}}$,
which requires
$c_{\mathrm{ghost}}=8k/(k+3)-3(2k+1)/(k+3)
=(8k-6k-3)/(k+3)=(2k-3)/(k+3)$.

For this to equal $4$: $2k-3=4(k+3)=4k+12$, giving
$-2k=15$, $k=-15/2$.  So $c_{\mathrm{ghost}}=4$ only at
$k=-15/2$, not for all~$k$.

The correct formula: $c_{\mathrm{ghost}}$ for subregular
DS of $\mathfrak{sl}_3$ is $(2k-3)/(k+3)$
(level-dependent, not a constant).  This is a key
difference from the principal case where
$c_{\mathrm{ghost}}=N(N-1)$ is level-independent.
The level-dependence arises because the subregular
embedding changes the conformal weights of the ghost
fields as a function of~$k$.
\end{computation}

\subsection{Modular extension: the ordered bar complex at
genus $g\ge 1$}
\label{subsec:ordered-bar-modular}

\begin{construction}[Modular ordered bar complex;
\ClaimStatusProvedHere]
\label{constr:modular-ordered-bar}
\index{bar complex!ordered!modular extension|textbf}
\index{modular R-matrix|textbf}
At genus~$g\ge 1$, the ordered configuration space on a curve
$\Sigma_g$ acquires new topology: the surface pure braid group
$\pi_1(\mathrm{Conf}_n^{\mathrm{ord}}(\Sigma_g))$ has generators
from the braid group (as at genus~$0$) and from the handle
loops of~$\Sigma_g$.  The monodromy of the ordered bar connection
around the handle loops produces genus corrections to the
$R$-matrix.

The \emph{modular ordered bar complex} is
\[
{\Barch}^{\mathrm{ord,mod}}(\cA)(g,n)
\;=\;
\bigoplus_{\Gamma\in\mathsf{Gr}^{\mathrm{st,ord}}_{g,n}}
\frac{\hbar^{g(\Gamma)}}{|\operatorname{Aut}(\Gamma)|}
\;(s^{-1}\bar\cA)^{\otimes|V(\Gamma)|}
\;\otimes\;
C_\bullet\!\bigl(
\mathrm{Conf}^{\mathrm{ord}}_{V(\Gamma)}(\Sigma_g)\bigr)
\;\otimes\;\det(E(\Gamma)),
\]
where $\mathsf{Gr}^{\mathrm{st,ord}}_{g,n}$ denotes stable
graphs whose vertices carry a \emph{planar} (total) ordering.

The \emph{modular $R$-matrix} is the genus expansion
\begin{equation}\label{eq:modular-r-matrix}
R^{\mathrm{mod}}(z;\hbar)
\;=\;
\sum_{g\ge 0}\hbar^{2g}\,r_g(z),
\end{equation}
with $r_0(z)=R(z)$ the genus-$0$ $R$-matrix and $r_g(z)$
($g\ge 1$) the genus-$g$ correction from the period structure
of~$\Sigma_g$.

At genus~$1$: the correction $r_1(z)$ arises from the monodromy
of the ordered bar connection around the $A$- and $B$-cycles
of the elliptic curve.  For
$\widehat{\mathfrak{sl}}_2$ at level~$k$:
\[
r_1(z)
\;=\;
\frac{\kappa}{(k+2)^2}\,
\frac{\wp(z;\tau)}{z}\,\Omega
\;+\;\text{(non-Casimir terms)},
\]
where $\wp$ is the Weierstrass function of the elliptic
curve with modulus~$\tau$, and $\kappa=3k/(k+2)$.
The Weierstrass function appears because the genus-$1$
propagator $\varpi_{12}=\partial_{z_1}\partial_{z_2}\log
\theta_1(z_1-z_2;\tau)$ has an expansion
$\varpi_{12}=1/(z_1-z_2)^2-\wp(z_1-z_2;\tau)+\cdots$
whose non-rational part is~$\wp$.  The genus-$0$
$R$-matrix $R(z)=1+\Omega/z+\cdots$ acquires an
elliptic deformation $R_\tau(z)=1+\Omega/z
+\hbar^2\kappa\,\wp(z;\tau)\,\Omega/(k+2)^2+\cdots$.

For $\cH_k$ (Heisenberg): $r_1(z)=k\,\wp(z;\tau)/z$
(scalar).  The modular $R$-matrix is
$R^{\mathrm{mod}}(z;\hbar)=z^k+\hbar^2 k\,\wp(z;\tau)/z+\cdots$:
the Heisenberg braiding acquires elliptic corrections at
genus~$1$.
\end{construction}

